\begin{figure}[t]
\centering
\begin{tikzpicture}[>=stealth, node distance=1.8cm, every node/.style={font=\small}]
    % Nodes
    \node[draw, minimum size=1cm, shape=rectangle] (S1) {$s_1$};
    \node[draw, minimum size=1cm, shape=rectangle, right=of S1] (S2) {$s_2$};
    \node[draw, minimum size=1cm, shape=rectangle, draw=none, right=of S2] (dots) {\dots};
    \node[draw, minimum size=1cm, shape=rectangle, right=of dots] (Snm1) {$s_{n-1}$};
    \node[draw, minimum size=1cm, shape=rectangle, right=of Snm1] (Sn) {$s_n$};

    % Self-loops
    \draw[->] (S1) edge[loop above] node[above] {$(s_1, c)$} coordinate[pos=0.5] (S1loop) (S1);
    \filldraw[black] (S1loop) circle (2pt);

    \draw[->] (S2) edge[loop above] node[above] {$(s_2, c)$} coordinate[pos=0.5] (S2loop) (S2);
    \filldraw[black] (S2loop) circle (2pt);

    \draw[->] (Snm1) edge[loop above] node[above] {$(s_{n-1}, c)$} coordinate[pos=0.5] (Snm1loop) (Snm1);
    \filldraw[black] (Snm1loop) circle (2pt);

    \draw[->] (Sn) edge[loop above] node[above] {$(s_n, c)$} coordinate[pos=0.5] (Snloop) (Sn);
    \filldraw[black] (Snloop) circle (2pt);

    % Transitions and splits
    \draw[-] (S1) -- node[below] {} ($(S1)!0.5!(S2)$);
    \filldraw[black] ($(S1)!0.5!(S2)$) circle (2pt) 
        node[label={90:$(s_1, b)$}] {};
    \draw[->] ($(S1)!0.5!(S2)$) -- (S2);
    \draw[->] ($(S1)!0.5!(S2)$) .. controls +(0,-1) and +(0,-1) .. (S1);

    \draw[-] (S2) -- node[below] {} ($(S2)!0.5!(dots)$);
    \filldraw[black] ($(S2)!0.5!(dots)$) circle (2pt) 
        node[label={90:$(s_2, b)$}] {};
    \draw[->] ($(S2)!0.5!(dots)$) -- (dots);
    \draw[->] ($(S2)!0.5!(dots)$) .. controls +(0,-1.5) and +(0,-1.5) .. (S1);

    \draw[-] (dots) -- node[below] {} ($(dots)!0.5!(Snm1)$);
    \filldraw[black] ($(dots)!0.5!(Snm1)$) circle (2pt)
        node[label={90:$(s_{n-2}, b)$}] {};
    \draw[->] ($(dots)!0.5!(Snm1)$) -- (Snm1);
    \draw[->] ($(dots)!0.5!(Snm1)$) .. controls +(0,-2) and +(0,-2) .. (S1);

    \draw[-] (Snm1) -- node[below] {} ($(Snm1)!0.5!(Sn)$);
    \filldraw[black] ($(Snm1)!0.5!(Sn)$) circle (2pt) 
        node[label={90:$(s_{n-1}, b)$}] {};
    \draw[->] ($(Snm1)!0.5!(Sn)$) -- (Sn);
    \draw[->] ($(Snm1)!0.5!(Sn)$) .. controls +(0,-2.5) and +(0,-2.5) .. (S1);

\end{tikzpicture}

\caption{A RNG that shows the worst-case complexity of the procedure for weak attractors. In the calculation of the weak attractor of $\set{s_0}$ for the square player. In the $k$-iteration $s_n-k$ is the only non-nature state added to $C_k$. Therefore $n$ iterations are needed. Whereas the computation of $B_k$ requires traversing the $n$ states of the upper chain, taking $O(|\Delta|)$ steps. }
\label{fig:example_worst_case}
\end{figure}